% Module 2: Rules of the Quantum World
% Estimated slides: 30 | Duration: ~2 hours
% Key topics: Superposition (revisited), measurement, entanglement, interference,
%             no-cloning theorem, decoherence, double-slit intuition, Bell inequality
% Labs: Entanglement demo in Qiskit; interference visualization

%%%%%%%%%%%%%%%%%%%%%%%%%%%%%%%%%%%%%%%%%%%%%%%%%%%%%%%%%%
\begin{frame}[fragile]\frametitle{}
\begin{center}
{\Large Module 2: Rules of the Quantum World}
\end{center}
\end{frame}

%%%%%%%%%%%%%%%%%%%%%%%%%%%%%%%%%%%%%%%%%%%%%%%%%%%%%%%%%%%
\begin{frame}[fragile]\frametitle{The Quantum Rule Book}
\begin{itemize}
\item Quantum mechanics has a small set of fundamental rules
\item These rules are counterintuitive but extremely well-tested experimentally
\item We will cover 6 rules that matter most for quantum computing:
  \begin{enumerate}
  \item Superposition
  \item Measurement and Collapse
  \item Entanglement
  \item Interference
  \item No-Cloning
  \item Decoherence
  \end{enumerate}
\item Think of these as the ``physics laws'' your programs must respect
\end{itemize}
\end{frame}

%%%%%%%%%%%%%%%%%%%%%%%%%%%%%%%%%%%%%%%%%%%%%%%%%%%%%%%%%%%
\begin{frame}[fragile]\frametitle{Rule 1: Superposition}
\begin{itemize}
\item A quantum system can exist in multiple states simultaneously
\item Only takes a definite value when measured
\item Not just a lack of knowledge about the state --- the system genuinely has no definite value
\item Demonstrated by the double-slit experiment: a single particle goes through both slits
\item For qubits: can be in both 0 and 1 at the same time
\item Enables quantum parallelism: apply operations to all states simultaneously
\end{itemize}
% Suggested diagram: Double-slit experiment schematic showing interference pattern
\end{frame}

%%%%%%%%%%%%%%%%%%%%%%%%%%%%%%%%%%%%%%%%%%%%%%%%%%%%%%%%%%%
\begin{frame}[fragile]\frametitle{The Double-Slit Experiment (Intuition)}
\begin{itemize}
\item Shine light (or electrons) through two slits onto a screen
\item Classical particle: goes through one slit, makes two bands on screen
\item Quantum particle: behaves like a wave, goes through both slits simultaneously
\item Result: interference pattern (alternating bright and dark bands)
\item If you try to detect which slit the particle went through: interference disappears!
\item Observation itself changes the quantum system
\item Moral: quantum objects don't have a definite path until observed
\end{itemize}
% Suggested diagram: Double-slit experiment diagram with classical vs quantum results
\end{frame}

%%%%%%%%%%%%%%%%%%%%%%%%%%%%%%%%%%%%%%%%%%%%%%%%%%%%%%%%%%%
\begin{frame}[fragile]\frametitle{Rule 2: Measurement and Collapse}
\begin{itemize}
\item Measuring a qubit in superposition gives a random 0 or 1 outcome
\item The probability of each outcome is set by the amplitudes
\item After measurement: the qubit is in a definite state (collapsed)
\item Measurement is \textit{irreversible} --- you cannot undo it
\item Quantum gates (before measurement) are reversible; measurement is not
\item Always measure at the end of a quantum circuit, not in the middle
\end{itemize}
% Suggested diagram: Circuit with gates, then measurement arrow, then classical bit output
\end{frame}

%%%%%%%%%%%%%%%%%%%%%%%%%%%%%%%%%%%%%%%%%%%%%%%%%%%%%%%%%%%
\begin{frame}[fragile]\frametitle{Rule 3: Entanglement}
\begin{itemize}
\item Two or more qubits can be entangled: their states are correlated
\item Measuring one entangled qubit instantly determines the other's measurement outcome
\item This correlation holds regardless of distance between the qubits
\item Einstein called it ``spooky action at a distance'' --- he didn't like it
\item Bell's experiments (1980s+) confirmed entanglement is real
\item Entanglement is not telepathy: you cannot use it to send information faster than light
\end{itemize}
% Suggested diagram: Two entangled qubits shown as linked nodes; measure one -> both collapse
\end{frame}

%%%%%%%%%%%%%%%%%%%%%%%%%%%%%%%%%%%%%%%%%%%%%%%%%%%%%%%%%%%
\begin{frame}[fragile]\frametitle{Entanglement: The Glove Analogy}
\begin{itemize}
\item Imagine you put one glove in each of two boxes, seal them, and send them to different cities
\item When one person opens their box and finds a left glove, they know the other has the right one
\item Classical: the gloves had definite left/right identities when packed --- no mystery
\item Quantum: the gloves had NO definite identities until one was opened --- truly random until then
\item Bell's theorem and experiments rule out the classical (pre-assigned) explanation
\item Entanglement is a real non-classical correlation
\end{itemize}
\end{frame}

%%%%%%%%%%%%%%%%%%%%%%%%%%%%%%%%%%%%%%%%%%%%%%%%%%%%%%%%%%%
\begin{frame}[fragile]\frametitle{Creating Entanglement: Bell States}
\begin{itemize}
\item The simplest entangled states are the 4 Bell states
\item $|\Phi^+\rangle = (|00\rangle + |11\rangle)/\sqrt{2}$: measure both qubits, always get same result
\item If qubit A is 0, qubit B is guaranteed to be 0; if A is 1, B is 1
\item Created by: Hadamard gate on qubit 0, then CNOT between qubit 0 and 1
\item This is one of the most important 2-qubit operations in quantum computing
\end{itemize}
\begin{verbatim}
qc = QuantumCircuit(2, 2)
qc.h(0)        # Superposition on qubit 0
qc.cx(0, 1)    # CNOT: entangle qubit 0 and 1
qc.measure([0,1], [0,1])
\end{verbatim}
% Suggested diagram: Circuit diagram for Bell state preparation
\end{frame}

%%%%%%%%%%%%%%%%%%%%%%%%%%%%%%%%%%%%%%%%%%%%%%%%%%%%%%%%%%%
\begin{frame}[fragile]\frametitle{Rule 4: Interference}
\begin{itemize}
\item Quantum amplitudes can add (constructive) or cancel (destructive)
\item This is similar to wave interference in optics and sound
\item Quantum algorithms use interference to amplify the probability of the right answer
\item Constructive interference: two paths to the correct answer add up
\item Destructive interference: paths to wrong answers cancel out
\item Interference is what makes quantum algorithms faster than brute force
\end{itemize}
% Suggested diagram: Wave diagram showing constructive and destructive interference; map to circuit
\end{frame}

%%%%%%%%%%%%%%%%%%%%%%%%%%%%%%%%%%%%%%%%%%%%%%%%%%%%%%%%%%%
\begin{frame}[fragile]\frametitle{Rule 5: No-Cloning Theorem}
\begin{itemize}
\item It is impossible to make a perfect copy of an unknown quantum state
\item Classical bits can be copied: $0 \to 00$, $1 \to 11$
\item Qubits cannot: copying would violate the laws of quantum mechanics
\item This is not a technology limitation --- it is provably impossible
\item \textbf{Why it matters for computing}: You cannot ``backup'' quantum data mid-algorithm
\item \textbf{Why it matters for security}: An eavesdropper cannot clone quantum encryption keys
\item This is the foundation of quantum cryptography
\end{itemize}
\end{frame}

%%%%%%%%%%%%%%%%%%%%%%%%%%%%%%%%%%%%%%%%%%%%%%%%%%%%%%%%%%%
\begin{frame}[fragile]\frametitle{Rule 6: Decoherence}
\begin{itemize}
\item Real qubits interact with their environment (temperature, electromagnetic noise)
\item These interactions cause qubits to lose their quantum properties over time
\item Decoherence: the qubit behaves increasingly like a classical bit with noise
\item T1 time: how long a qubit stays in excited state $|1\rangle$ (energy relaxation)
\item T2 time: how long phase information survives (dephasing)
\item Typical IBM superconducting qubits: T1, T2 in microseconds to milliseconds
\item This limits how deep (long) our circuits can be
\end{itemize}
% Suggested diagram: Qubit coherence curve (amplitude vs time decaying)
\end{frame}

%%%%%%%%%%%%%%%%%%%%%%%%%%%%%%%%%%%%%%%%%%%%%%%%%%%%%%%%%%%
\begin{frame}[fragile]\frametitle{Schrodinger's Cat: A Famous Thought Experiment}
\begin{itemize}
\item A cat is in a sealed box with a quantum device that may or may not release poison
\item Before opening: the cat is in a superposition of alive AND dead
\item After opening (measurement): the cat is definitely one or the other
\item Intended to highlight the absurdity of applying quantum rules to macroscopic objects
\item Reality: large objects decohere almost instantly --- superpositions don't survive
\item Quantum superpositions only survive at the scale of atoms, electrons, photons
\end{itemize}
% Suggested diagram: Schrodinger's cat box with quantum/classical boundary
\end{frame}

%%%%%%%%%%%%%%%%%%%%%%%%%%%%%%%%%%%%%%%%%%%%%%%%%%%%%%%%%%%
\begin{frame}[fragile]\frametitle{Bell's Theorem: Quantum vs Classical Correlations}
\begin{itemize}
\item Bell (1964) devised a test to distinguish quantum correlations from classical ones
\item Classical correlations (hidden variable theory): limited by Bell inequalities
\item Quantum entanglement violates Bell inequalities
\item Experiments by Aspect (1982), Zeilinger (2022 Nobel) confirmed quantum violation
\item Conclusion: quantum mechanics is fundamentally different from classical probability
\item Practical consequence: entanglement enables protocols impossible classically (teleportation, QKD)
\end{itemize}
% Suggested diagram: CHSH inequality graph showing classical bound vs quantum violation
\end{frame}

%%%%%%%%%%%%%%%%%%%%%%%%%%%%%%%%%%%%%%%%%%%%%%%%%%%%%%%%%%%
\begin{frame}[fragile]\frametitle{Quantum Tunneling (Bonus: Not in Circuits, but Everywhere)}
\begin{itemize}
\item A quantum particle can pass through a barrier it classically couldn't cross
\item Analogy: a ball rolling up a hill, but it sometimes ``teleports'' to the other side
\item Enables: nuclear fusion in stars, scanning tunneling microscopes, flash memory
\item In quantum computing: used in quantum annealing (D-Wave) for optimization
\item Not directly used in gate-model quantum computing (our course focus)
\item Shows how quantum effects underlie everyday technology
\end{itemize}
% Suggested diagram: Classical ball vs quantum particle at a potential barrier
\end{frame}

%%%%%%%%%%%%%%%%%%%%%%%%%%%%%%%%%%%%%%%%%%%%%%%%%%%%%%%%%%%
\begin{frame}[fragile]\frametitle{Quantum vs Classical: Side-by-Side Comparison}
\begin{itemize}
\item \textbf{State}: Classical = definite, Quantum = superposition until measured
\item \textbf{Copying}: Classical = free, Quantum = forbidden (no-cloning)
\item \textbf{Measurement}: Classical = passive reading, Quantum = disturbs and collapses
\item \textbf{Correlation}: Classical = can be explained locally, Quantum = entanglement is non-local
\item \textbf{Interference}: Classical = none, Quantum = amplitudes add/cancel
\item \textbf{Error}: Classical = bit flip is detectable, Quantum = also phase flip errors
\end{itemize}
% Suggested diagram: Comparison table side-by-side
\end{frame}

%%%%%%%%%%%%%%%%%%%%%%%%%%%%%%%%%%%%%%%%%%%%%%%%%%%%%%%%%%%
\begin{frame}[fragile]\frametitle{Lab: Entanglement and Measurement Correlation}
\begin{itemize}
\item Create a Bell state circuit in Qiskit
\item Run 1000 times on simulator
\item Observe: only 00 and 11 outcomes (never 01 or 10)
\item This demonstrates perfect correlation of entangled qubits
\item Extension: add noise and see how correlations degrade
\end{itemize}
\begin{verbatim}
from qiskit_aer import AerSimulator
from qiskit.visualization import plot_histogram
sim = AerSimulator()
# Run Bell state circuit and plot
\end{verbatim}
\end{frame}

%%%%%%%%%%%%%%%%%%%%%%%%%%%%%%%%%%%%%%%%%%%%%%%%%%%%%%%%%%%
\begin{frame}[fragile]\frametitle{Lab: Interference Demonstration}
\begin{itemize}
\item Create a circuit: H gate then H gate again (H applied twice = identity)
\item This is constructive and destructive interference in action
\item H$|0\rangle = |+\rangle$, then H$|+\rangle = |0\rangle$ (interference restores original state)
\item Demonstrates that phases matter: $|+\rangle$ and $|-\rangle$ look the same in Z measurements
\item But when interfered, the phase difference creates a measurable difference
\end{itemize}
\begin{verbatim}
qc = QuantumCircuit(1, 1)
qc.h(0)     # Create superposition
qc.h(0)     # Interfere back to |0>
qc.measure(0, 0)  # Should always give 0
\end{verbatim}
% Suggested diagram: HH circuit; ball analogy for constructive interference
\end{frame}

%%%%%%%%%%%%%%%%%%%%%%%%%%%%%%%%%%%%%%%%%%%%%%%%%%%%%%%%%%%
\begin{frame}[fragile]\frametitle{Exercise: Quantum Rule Identification}
\begin{itemize}
\item Identify which quantum rule is at play in each scenario:
\begin{enumerate}
  \item A qubit is in state $|+\rangle$; after measurement it always gives 0 in future -- (Collapse)
  \item Two qubits always agree when measured, even when far apart -- (Entanglement)
  \item An algorithm boosts probability of the correct answer using phase tricks -- (Interference)
  \item After 100 microseconds, a qubit's quantum state degrades -- (Decoherence)
  \item You cannot duplicate an unknown qubit to back it up -- (No-Cloning)
\end{enumerate}
\end{itemize}
\end{frame}

%%%%%%%%%%%%%%%%%%%%%%%%%%%%%%%%%%%%%%%%%%%%%%%%%%%%%%%%%%%
\begin{frame}[fragile]\frametitle{Module 2 Summary}
\begin{itemize}
\item \textbf{Superposition}: qubits can be in multiple states simultaneously
\item \textbf{Measurement}: collapses superposition; irreversible
\item \textbf{Entanglement}: non-local correlations between qubits; experimentally confirmed
\item \textbf{Interference}: amplitudes add/cancel; enables algorithmic speedup
\item \textbf{No-cloning}: unknown quantum states cannot be copied; basis of quantum crypto
\item \textbf{Decoherence}: quantum properties decay due to environmental interaction; limits circuit depth
\item Next: how to control qubits with quantum gates and build circuits
\end{itemize}
\end{frame}
